% Library CV waste_fr — transcription of writing-samples/cvs/Moumen_Nermine_6.pdf (CV09, 2026-09-25)
\documentclass[10pt]{article}
\usepackage[letter]{nermine-cv}
\begin{document}
\cvheader{Moumen Nermine}{Étudiante en dernière année du cycle ingénieur en Chimie Industrielle}{nermine.moumen@insat.ucar.tn}{+21627040414}{linkedin.com/in/nermine-moumen}

\section{FORMATION}
\entry{Institut National des Sciences Appliquées et de Technologie (INSAT)}{2022–Présent}
\summary{Étudiante en Chimie Industrielle, spécialité Procédés et Environnement, avec une formation solide en traitement et valorisation des déchets, complétée par des bases en génie des procédés, chimie analytique et HSE en laboratoire. Orientée vers l'économie circulaire et l'optimisation des filières de traitement des déchets. À la recherche d'un PFE.}

\section{EXPÉRIENCE PROFESSIONNELLE}
\entry{SEGULA Technologies - Stage en Modélisation Électrochimique des batteries}{06/2026–08/2026}
\begin{bul}
\item Développement d'un modèle numérique (Single Particle Model, MATLAB) d'une cellule Li-ion, résolu par FVM.
\item Optimisation du maillage et des coefficients de diffusion via une analyse de données, réduisant l'erreur RMSE de 80 \%.
\item Validation des résultats par une approche itérative et une analyse critique avec les données expérimentales et COMSOL.
\end{bul}
\entry{Control Union - Stage en Contrôle Qualité}{08/2025–09/2025}
\begin{bul}
\item Analyse de lubrifiants (viscosité et métaux lourds par \textbf{ICP-OES}).
\item Préparation et interprétation de rapports d'analyse pour vérifier la conformité aux normes \textbf{ASTM D5185} et \textbf{ISO 17025}.
\end{bul}
\entry{ETAP - Stage en Caractérisation d'hydrocarbures et traitement des Eaux industrielles}{07/2024–08/2024}
\begin{bul}
\item Formation aux techniques de géochimie organique (GC, GC-MS, pyrolyse Rock-Eval) pour la caractérisation d'hydrocarbures.
\item Analyses physico-chimiques des eaux (pH, salinité, bicarbonates, sulfates, chlorures) par \textbf{chromatographie ionique} et \textbf{ICP-AES}.
\item Étude des méthodes de traitement des eaux pétrolières par échange d'ions et par coagulation-floculation.
\end{bul}
\section{EXPÉRIENCE COMPLÉMENTAIRE}
\entry{Cascade - Enseignante de Chimie en Ligne (Temps partiel)}{2024–Présent}
\begin{bul}
\item Accompagnement individualisé en chimie pour des élèves de Terminale française, couvrant les concepts du niveau examen.
\end{bul}

\section{PROJETS ACADÉMIQUES}
\project{Valorisation énergétique des déchets organiques par méthanisation (ASBR)}
\begin{bul}
\item Étude de la chaîne de traitement des déchets organiques : réception, tri, broyage, homogénéisation et préparation de la charge.
\item Caractérisation et suivi expérimental de \textbf{5 configurations de traitement} (boues primaires, secondaires, mixtes et co-digestion avec déchets fruits/légumes) dans des digesteurs ASBR (TRH = 20 j).
\item Analyse des performances des différentes filières à travers la DCO, MVS, charge volumique, production de biogaz et teneur en CH$_4$, avec jusqu'à \textbf{0,50 L/j de biogaz à 64,52 \% CH$_4$}.
\end{bul}
\project{Gestion et traitement des déchets radioactifs : stratégies de confinement, de stabilisation et de valorisation}
\begin{bul}
\item Étude comparative des filières de gestion des déchets nucléaires : entreposage en piscine et entreposage à sec en surface.
\item Analyse du procédé d'immobilisation par vitrification et du stockage géologique profond pour le confinement des déchets.
\item Étude des voies de réduction de la radiotoxicité par transmutation nucléaire et par les réacteurs nucléaires de 4\textsuperscript{e} génération.
\end{bul}
\project{Analyse techno-environnementale de la production d'éthylène : vapocraquage vs déshydratation du bioéthanol}
\begin{bul}
\item Simulation comparative de deux voies de production sous \textbf{Aspen Plus V14}, avec établissement des bilans matière et énergie et des diagrammes de flux (PFD) du procédé, pour évaluer la faisabilité technique de chaque filière.
\item Analyse de Cycle de Vie cradle-to-gate sous \textbf{OpenLCA} pour comparer les impacts environnementaux des deux voies.
\end{bul}

\section{COMPÉTENCES TECHNIQUES}
\begin{bul}
\item \textbf{Techniques de Laboratoire :} Chromatographie (HPLC, GC, CCM), Spectrométrie (MS, ICP-OES), Spectroscopie (RMN, FTIR, UV-Vis), Diffraction (XRD), Microscopie (MEB, MET), Électrochimie (Voltamétrie cyclique, polarographie).
\item \textbf{Simulation :} Aspen Plus et HYSYS, COMSOL Multiphysics\sep\textbf{Programmation :} \textit{MATLAB, Mathematica, Python}
\item \textbf{Données \& Bureautique :} Pack Microsoft Office, OpenLCA\sep\textbf{Normes :} ISO 9001, ISO 14001, ISO 45001, ISO 17025, ISO 14040
\end{bul}

\section{ACTIVITÉS EXTRASCOLAIRES}
\textbf{IEEE INSAT} — Chef d'organisation de TRSYP 3.0 ; Responsable des activités sociales – MENASYP ; Présidente de l'EMBS.\\
\textbf{THEATRO INSAT} — Membre de l'équipe de danse.\par

\section{LANGUES ET CERTIFICATS}
\twocol{Langues}{Anglais (Avancé), Français (Bilingue), Arabe (Natif).}{Certificats}{Aspen HYSYS – MATLAB}
\end{document}
